\documentclass[8pt]{article}
\usepackage{graphicx} % Required for inserting images
\usepackage{amsthm}
\usepackage{amsmath}
\usepackage{amsfonts}
\usepackage{amsopn}
\usepackage{comment}
\usepackage{amssymb}
\usepackage{hyperref}
\usepackage{bussproofs}
\usepackage[all, 2cell]{xy}
\usepackage[all]{xy}
\usepackage{rotating}
\usepackage{lscape}
\usepackage{minted}
\usepackage{dsfont}
\usepackage{tikz}
\usepackage{tikz-cd}
\usepackage{stmaryrd}
\usepackage{multicol}
\usepackage{cmll}
\usepackage{wasysym}
\usetikzlibrary{shapes.geometric, arrows, positioning, backgrounds}
\usetikzlibrary{cd}
\makeatletter
\tikzcdset{
  eq node/.style={
    commutative diagrams/math mode=false, anchor=center},
  eq/.style={
    phantom,
    /tikz/every to/.append style={
      edge node={node[commutative diagrams/eq node]
        {\@eqnswtrue\make@display@tag\ltx@label{#1}}}}}}
\makeatother

\theoremstyle{definition}
\newtheorem{definition}{Definition}[section]

\theoremstyle{definition}
\newtheorem{theorem}{Theorem}[section]

\theoremstyle{definition}
\newtheorem{claim}{Claim}[section]

\theoremstyle{definition}
\newtheorem{ex}{Example}[section] 

\theoremstyle{definition}
\newtheorem{cons}{Construction}[section] 

\theoremstyle{definition}
\newtheorem{rem}{Remark}[section] 


\theoremstyle{definition}
\newtheorem{prop}{Proposition}[section]

\theoremstyle{definition}
\newtheorem{lemma}{Lemma}[section]

\theoremstyle{definition}
\newtheorem{fact}{Fact}[section]

\theoremstyle{definition}
\newtheorem{remark}{Remark}[section]

\theoremstyle{definition}
\newtheorem{notation}{Notation}[section]

\theoremstyle{definition}
\newtheorem{example}{Example}[section]

\theoremstyle{definition}
\newtheorem{exercise}{Exercise}[section]

\theoremstyle{definition}
\newtheorem{col}{Corollary}[section]

\theoremstyle{question}
\newtheorem{question}{Question}

\let\strokeL\L
\renewcommand\L{\mathbf{L}}

\DeclareMathOperator{\bang}{!}

\DeclareMathOperator*{\colim}{\operatorname{colim}}
\newcommand{\Ob}[1]{\operatorname{Ob}({\mathcal{#1}})}
\newcommand{\Cov}[1]{\operatorname{Cov}(#1)}

\newcommand{\circvee}{\mathbin{\text{\ooalign{\hfil$\vee$\hfil\cr\kern0.1ex$\bigcirc$}}}}
\newcommand{\circwedge}{\mathbin{\text{\ooalign{\hfil$\wedge$\hfil\cr\kern0.1ex$\bigcirc$}}}}

\title{Braided Notions of Dialogue Categories: reworked proofs}
\author{Daniel Rogozin \footnote{The original text is by Paul-Andr\'{e} Melliès}}
\date{ }

\newcommand{\IntHom}[1]{[#1]}

\begin{document}

\maketitle

\section{Dialogue Categories: Basic Constructions and Facts}

\begin{definition}
  Let $\mathcal{C}$ be a monoidal category, a $\mathcal{C}$ is \emph{dialogue category}
  if there is an exponentiable object $\bot$, i.e., there are bijections natural in $A$ and $B$
  \begin{center}
    $\varphi_{A,B} : \mathcal{C}(A \otimes B, \bot) \cong (B, A \multimap \bot)$

    $\psi_{A, B} : \mathcal{C}(A \otimes B, \bot) \cong (A, \bot \multimapinv B )$
  \end{center}
\end{definition}

Here are some useful combinators:
\begin{itemize}
  \item The left evaluation map:
  \[
  \operatorname{leval}_A = \varphi^{-1}_{A,A \multimap \bot} 1_{A \multimap \bot} : A \otimes (A \multimap \bot) \to \bot
  \]
  \item The right evaluation map:
  \[
  \operatorname{reval}_A = \psi^{-1}_{\bot \multimapinv A, A} 1_{\bot \multimapinv A} : (\bot \multimapinv A) \otimes A \to \bot
  \]
  \item The composite left evaluation map:
  \begin{center}
  $ \operatorname{lev}_{B,A} : B \otimes ((A \otimes B) \multimap \bot) \to (A \multimap \bot)$ \\
  $\operatorname{lev}_{B,A} = \varphi_{A, B \otimes ((A \otimes B) \multimap \bot)} (\operatorname{leval}_{A \otimes B} \circ \alpha_{A, B, A \otimes B \multimap \bot})$
  \end{center}
  \item The composite right evaluation map:
  \begin{center}
  $\operatorname{rev}_{A, B} : (\bot \multimapinv (A \otimes B)) \otimes A \to (\bot \multimapinv B)$ \\
  $\operatorname{rev}_{A, B} = \psi_{(A \otimes B \multimapinv \bot) \otimes A, B} (\operatorname{reval}_{A \otimes B} \circ \alpha_{A \otimes B \multimapinv \bot, A, B} )$
  \end{center}
\end{itemize}

\begin{definition}
   Let $\mathcal{C}$ be a dialogue category,
  \begin{itemize} 
  \item a \emph{turn} in $\mathcal{C}$ is a natural isomorphism
  $\operatorname{turn}_A : A \multimap \bot \cong \bot \multimapinv A$.
  \item a \emph{wheel} in $\mathcal{C}$ is a bijection $\mathcal{C}(A \otimes B, \bot) \simeq \mathcal{C}(B \otimes A, \bot)$
  natural in both $A$ and $B$. 
  \end{itemize}
\end{definition}

\begin{theorem}
  There is a 1-to-1 correspondence between turns and wheels in any dialogue category $\mathcal{C}$.
\end{theorem}
\begin{proof}
  Given a turn $\operatorname{turn}_A : A \multimap \bot \cong \bot \multimapinv A$ natural in $A$,
  define $\operatorname{wheel}_{A,B} : \mathcal{C}(A \otimes B, \bot) \to \mathcal{C}(B \otimes A, \bot)$
  as the following composite of bijections:
  \[
  \mathcal{C}(A \otimes B, \bot) \xrightarrow{\varphi_{A,B}} \mathcal{C}(B, A \multimap \bot) \xrightarrow{\mathcal{C}(1_B, \operatorname{turn}_A)} \mathcal{C}(B, \bot \multimapinv A) \xrightarrow{\psi^{-1}_{A,B}} \mathcal{C}(B \otimes A, \bot)
  \]
  This is obviously invertible, just take $\varphi^{-1}_{A,B} \circ \mathcal{C}(B, A \multimap \bot) \circ \psi_{A,B}$.
  $\operatorname{wheel}$ is also natural in both arguments. 
  Take $f : A_1 \to A_2$ and $g : A_2 \otimes B \to \bot$, then the following square commutes:
  \[
  \begin{tikzcd}
    B \otimes A_1 \ar[rr, "1_B \otimes f"] \ar[dr, "\operatorname{wheel}(g \circ f \otimes 1_B)"'] && B \otimes A_2 \ar[dl, "\operatorname{wheel}_{A_2,B}(f)"] \\
    & \bot
  \end{tikzcd}
  \]

  First of all, we have the following by the naturality of internal homs:
  \begin{equation}~\label{eq:1}
  \operatorname{turn}_{A_1} \circ \: \mathcal{C}(f,\bot) = \mathcal{C}(f,\bot) \circ \operatorname{turn}_{A_2}
  \end{equation}
  Further, observe that:
  \begin{equation}~\label{e2:1}
    \psi^{-1}_{B,A_1}(\operatorname{turn}_{A_1} \circ \: \varphi_{A_1,B}(g \circ f \otimes 1_B)) = \psi^{-1}_{B,A_2} (\operatorname{turn}_{A_2} \circ \: \varphi_{A_2,B}(g)) \circ (1_B \otimes f)
  \end{equation}
  Since
  \begin{align*}
    \psi^{-1}_{B,A_1}(\operatorname{turn}_{A_1} \circ \: \varphi_{A_1,B}(g \circ f \otimes 1_B)) &= \psi^{-1}_{B,A_1}(\operatorname{turn}_{A_1} \circ \: \mathcal{C}(f,1_{\bot}) \circ \varphi_{A_2,B}(g)) \\
    &= \psi^{-1}_{B,A_1}(\mathcal{C}(f,1_{\bot}) \circ \operatorname{turn}_{A_2} \circ \varphi_{A_2,B}(g))
    &= \psi^{-1}_{B,A_2} (\operatorname{turn}_{A_2} \circ \: \varphi_{A_2,B}(g)) \circ (1_B \otimes f)
  \end{align*}
\end{proof}

\end{document}